\documentclass[11pt]{article}

\usepackage[margin=1in]{geometry}
\usepackage{times}
\usepackage{amsmath,amssymb,amsthm}
\usepackage{algorithm}
\usepackage{algorithmic}
\usepackage{enumitem}
\usepackage{graphicx}
\usepackage[hidelinks]{hyperref}
\usepackage[numbers]{natbib}
\usepackage{booktabs}
\usepackage{multirow}
\usepackage[table]{xcolor}
\usepackage{caption}
\usepackage{subcaption}
\usepackage{array}
\usepackage{tabularx}

\theoremstyle{definition}
\newtheorem{definition}{Definition}[section]

\newcommand{\Human}{\textsf{Human}}
\newcommand{\AgentA}{\textsf{AgentA}}
\newcommand{\AgentB}{\textsf{AgentB}}
\newcommand{\Shared}{\textit{shared}}

\title{Trust Calibration and Collaborative Decision-Making in\\
       Human--Agent Sequential Graph Colouring:\\
       A Within-Subjects Experimental Study}
\author{}
\date{}

\begin{document}
\maketitle

%% ============================================================
\begin{abstract}
%% ============================================================

We present a controlled experiment investigating how humans calibrate trust in AI agents
across differing levels of agent autonomy and plan quality.  The task is a sequential
graph-colouring problem in which a human participant and two rule-based AI agents jointly
colour the nodes of a graph, earning individual and shared scores that depend on colour
choices.  The payoff structure creates a deliberate tension: each participant's individually
preferred colour differs from the shared-optimal colour, so maximising the team score
requires coordination and a degree of individual sacrifice.  Three experimental conditions
vary the degree of agent autonomy and plan quality: in Mode~1, agents issue per-node
proposals with explicit rationales and the human decides; in Mode~2A, agents pre-plan the
full colouring with 80\% probability of picking the shared-optimal colour; in Mode~2B, the
same autonomous planning process is used but with only a 30\% probability of choosing the
shared-optimal colour.  Critically, the interface for Mode~2A and Mode~2B is visually
identical, so participants cannot detect quality differences from presentation alone.  We
measure shared score, individual score, plan-modification rate, explanation engagement, and
decision latency, together with score improvement across repeated attempts.  The study
offers a controlled platform for examining trust calibration, appropriate reliance, and the
interaction between agent transparency and agent competence in human--AI collaborative
settings.

\end{abstract}

%% ============================================================
\section{Introduction}
\label{sec:intro}
%% ============================================================

A central challenge in human--AI collaboration is calibrating trust appropriately: humans
should rely on AI agents when those agents are competent and override them when they are
not~\citep{lee2004trust,parasuraman2010complacency}.  Miscalibrated trust takes two forms:
\emph{overtrust}, in which users accept AI recommendations even when the AI is wrong, and
\emph{undertrust}, in which users override good recommendations unnecessarily.  Both forms
are costly, but overtrust receives more attention in safety-critical domains because it can
cause humans to defer to confident-looking but incorrect AI outputs.

A key variable shaping trust calibration is the \emph{transparency} of the AI's reasoning
process.  When an agent shows its rationale---why it recommends a particular action, what
information it used, what trade-offs it considered---users have more material with which to
evaluate reliability.  Autonomous agents that produce pre-made plans and present them for
human approval represent a particular transparency challenge: the user can see the
recommendation but not the planning process that produced it.

We study these dynamics using a graph-colouring task that operationalises individual-vs-team
trade-offs in a simple, reproducible setting.  The graph is divided among three participants
(one human, two AI agents), each controlling a distinct cluster of nodes.  Participants
colour their nodes sequentially, one at a time, in a fixed order.  The payoff structure is
designed so that each participant's selfishly preferred colour differs from the colour that
maximises the shared team score.  This creates a coordination dilemma: a participant who
always chooses their personally preferred colour earns more individually but less as a team;
a participant who always chooses the shared-optimal colour sacrifices personal gain for
collective benefit.

Three experimental conditions differ in how AI agent advice is delivered:

\begin{description}[leftmargin=2em,itemsep=2pt]
  \item[Mode~1 (Manual Collaborative)]  For each node in the human's cluster, both AI
    agents issue a proposed colour with a short written rationale and a table of expected
    scores.  The human sees these proposals simultaneously and then makes their own colour
    choice, unconstrained by the proposals.  Agent nodes are coloured automatically by
    the respective agent's own rule-based logic.
  \item[Mode~2A (Autonomous, High Quality)]  Before colouring begins, a planning agent
    runs a deliberation process over the full node sequence and produces a complete
    proposed colouring.  The plan is then presented to the human one step at a time; the
    human can accept or override each step.  The planning process picks the shared-optimal
    colour with probability 0.8, so roughly 80\% of plan steps are objectively good for
    the team.
  \item[Mode~2B (Autonomous, Low Quality)]  Structurally identical to Mode~2A.  The
    interface, terminology, and presentation are the same.  The only difference is that
    the planning process picks the shared-optimal colour with probability 0.3, producing
    substantially lower-quality plans.  Participants are not told which quality condition
    they are in.
\end{description}

This design allows us to ask: when the interface is held constant and only the underlying
plan quality changes, do participants adjust their override behaviour appropriately?  Do
participants who engage with agent rationales (by clicking to expand them) show better
calibration?  Does score improve across repeated attempts (up to three per session) as
participants learn the task and the agents' tendencies?

The rest of this paper is organised as follows.  Section~\ref{sec:related} reviews related
work.  Section~\ref{sec:task} defines the sequential graph-colouring task and its scoring
model.  Section~\ref{sec:graphs} describes the graph topologies used in the study.
Section~\ref{sec:conditions} specifies the three experimental conditions in detail.
Section~\ref{sec:agents} describes the agent models underlying each condition.
Section~\ref{sec:interface} describes the interface.
Section~\ref{sec:procedure} describes the experimental procedure.
Section~\ref{sec:measures} defines the dependent variables and hypotheses.
Section~\ref{sec:results} presents the results (to be completed).

%% ============================================================
\section{Related Work}
\label{sec:related}
%% ============================================================

\subsection{Trust in Automation and AI}

Lee and See~\citeyearpar{lee2004trust} define trust in automation as ``the attitude that an
agent will help achieve an individual's goals in a situation characterised by uncertainty
and vulnerability.''  A large body of literature has documented that trust does not
automatically track competence: humans exhibit automation bias (uncritical acceptance of
automated outputs)~\citep{parasuraman2010complacency} as well as automation disuse
(ignoring reliable automation).  Appropriate reliance is the goal: Buccinca et al.\
\citeyearpar{buccinca2021trust} show that cognitive forcing functions---interventions that
require users to reason before viewing AI recommendations---can reduce overtrust.  Our
design takes a different approach, varying the quality of the underlying agent while
holding the interface constant, to test whether participants detect quality differences
behaviourally.

\subsection{Human--AI Teaming in Optimisation Tasks}

Several studies have used combinatorial optimisation tasks to investigate human--AI
collaboration.  Wilder et al.~\citeyearpar{wilder2021learning} study human--AI teams on
scheduling problems, finding that humans override AI recommendations more when AI
explanations are present but not always more accurately.  Lai and Tan~\citeyearpar{lai2019human}
find that providing model explanations improves human decision accuracy only when
explanations correctly reflect the model's actual reasoning.  Our setting isolates a
similar question: does having access to an explanation (by clicking ``Explain'' in Mode~2)
correlate with better override decisions?

\subsection{Sequential Decision-Making and Path Dependency}

The sequential colouring structure of our task introduces path dependency: a suboptimal
choice early in the sequence constrains later options.  This mirrors real-world scenarios
such as project planning, resource allocation, and shift scheduling, where decisions
accumulate.  Kahneman et al.~\citeyearpar{kahneman1982judgment} document how early choices
anchor subsequent ones; in our setting, the first few node colour decisions may set an
inertial trajectory that the human then accepts rather than revises.  The repeated-attempt
structure (up to three attempts per session) allows us to measure whether and how quickly
participants learn to break this inertia.

\subsection{Graph Colouring as a Coordination Task}

Graph colouring has been used as a canonical coordination task in distributed constraint
optimisation~\citep{modi2005adopt,rogers2011bounded}.  Unlike classical DCOP settings,
our version does not penalise constraint violations as hard constraints---all colourings
are valid, but only some are \emph{optimal} under a shared scoring function.  This makes
the task an \emph{optimisation under conflicting incentives} rather than a feasibility
problem, which is both more realistic (agents' interests often misalign partially rather
than completely) and allows continuous measurement of coordination quality via score.

%% ============================================================
\section{Task: Sequential Graph Colouring}
\label{sec:task}
%% ============================================================

\subsection{Graph Structure}

The experimental task uses an undirected graph $G = (V, E)$ partitioned among three
participants: a human $h$ and two AI agents $a_1$ (\AgentA) and $a_2$ (\AgentB).  Each
participant controls a disjoint cluster of nodes:

\[
  V = V_h \sqcup V_{a_1} \sqcup V_{a_2}
\]

where $V_p$ denotes the set of nodes owned by participant $p$.  Edges $E$ include
both \emph{intra-cluster} edges (between nodes in the same cluster) and
\emph{inter-cluster} edges (between nodes in different clusters).  Inter-cluster edges
are the coordination interface: a colour chosen for a node in one cluster may clash with
or constrain choices in adjacent clusters.

Each participant assigns a colour from the shared domain $\mathcal{C} = \{\text{red},
\text{blue}, \text{green}\}$ to each of their nodes.  Unlike standard graph-colouring
problems, there is no \emph{hard} constraint requiring adjacent nodes to have different
colours.  Instead, adjacency conflicts incur a \emph{score penalty} described in
Section~\ref{sec:scoring}.  This design choice ensures that all colourings are valid (the
task always terminates) while retaining the incentive to avoid conflicts.

\subsection{Sequential Colouring}
\label{sec:sequential}

Nodes are coloured \emph{one at a time} in a fixed predetermined order.  The default
order is: all human-owned nodes ($V_h$), followed by all \AgentA-owned nodes
($V_{a_1}$), followed by all \AgentB-owned nodes ($V_{a_2}$).  Nodes within each
cluster follow the order specified by the graph definition.

This sequential structure introduces \emph{path dependency}: once a node is coloured,
its choice constrains (or at minimum influences) the optimal choices for subsequently
coloured adjacent nodes.  Because human nodes are coloured first, the human's decisions
constrain the agents more than the agents' decisions constrain the human.  This is an
intentional asymmetry: we are interested in how humans make decisions while anticipating
(or failing to anticipate) their downstream effects on agent scoring.

An optional \emph{random turn order} mode shuffles the ordering within each attempt.
When enabled with the \texttt{include\_agent\_turns} flag, agent nodes may be
interspersed with human nodes; when disabled (default), only human nodes are reordered
among themselves, and all agent nodes follow.  In both cases, the planning agent
(Mode~2) re-plans according to the actual shuffled order, ensuring the plan corresponds
to the sequence presented to the user.

Participants cannot skip or reorder nodes.  They colour nodes strictly in the presented
sequence: in Mode~1, they see agent proposals for the current node, then choose; in
Mode~2, they see the agent's pre-planned recommendation for the current node, then
accept or override.

\subsection{Scoring Model}
\label{sec:scoring}

The scoring system supports both individual and collective performance measurement.
Each colour assignment earns points for \emph{every} participant simultaneously, with
different amounts per participant.  Table~\ref{tab:default_points} shows the default
per-node point values.

\begin{table}[h]
\centering
\caption{Default per-node point values for each colour and participant.
         Shared total is the sum across all three participants.
         Each participant's individually preferred colour (maximum personal gain) is shown in bold.}
\label{tab:default_points}
\begin{tabular}{lcccc}
\toprule
\textbf{Colour} & \textbf{Human} & \textbf{AgentA} & \textbf{AgentB} & \textbf{Shared Total} \\
\midrule
\textbf{Red}   & \textbf{4}    & 1              & 2               & 7 \\
Blue            & 1             & \textbf{4}     & 2               & 7 \\
Green           & 2             & 2              & \textbf{4}      & 8 \\
\bottomrule
\end{tabular}
\end{table}

The key feature of this payoff structure is that:
\begin{enumerate}[nosep]
  \item Each participant has a \emph{different} individually preferred colour (the colour
    that maximises their personal points on any given node).
  \item The \emph{shared-optimal} colour --- the colour that maximises the sum of all
    participants' points --- is green (8 pts total), which is \AgentB's preferred colour
    but not the preferred choice for \Human{} or \AgentA.
  \item Choosing one's own preferred colour sacrifices 1 point of shared score per node
    relative to always choosing green.
\end{enumerate}

This creates a \emph{social dilemma} structure embedded within a sequential colouring
task.  A human who always chooses red earns the maximum personal score (4 pts/node)
but contributes only 7 pts/node to the shared pool.  A human who defers to green earns
only 2 pts/node personally but contributes 8 pts/node to the shared total.  Whether
participants sacrifice personal gain for shared benefit --- and whether agent advice
moves them in either direction --- is a primary outcome of interest.

\paragraph{Clash Penalty.}
If two adjacent nodes (sharing an edge in $E$) are assigned the same colour, a penalty
of 10~points is deducted from the shared score.  For intra-participant clashes (both
nodes owned by the same participant), the full penalty is deducted from that participant's
individual score.  For inter-participant clashes (adjacent nodes owned by different
participants), the penalty is split: $\lfloor 10/2 \rfloor = 5$ points from each owner.
Formally, the score for a complete assignment $x : V \to \mathcal{C}$ is:

\begin{equation}
  \text{Score}_p(x) = \sum_{v \in V_p} \text{pts}_p(x_v)
    - \sum_{\substack{(u,v) \in E \\ x_u = x_v}} \text{penalty}(u, v, p)
  \label{eq:score}
\end{equation}

where $\text{pts}_p(c)$ is participant $p$'s point value for colour $c$ (from
Table~\ref{tab:default_points}), and $\text{penalty}(u, v, p)$ is the portion of the
clash penalty on edge $(u,v)$ allocated to participant $p$.

The shared score is $\text{Score}_{\Shared}(x) = \sum_p \text{Score}_p(x)$.  The
maximally achievable shared score, with no clashes, is $8 \times |V|$ (all nodes
green).  The minimum achievable shared score depends on the graph topology and the
number of edges.

\paragraph{Score interpretation.}
Because the graph is fixed within a session, the shared score is bounded above and
below by constants determined by the topology.  The key observable is therefore the
\emph{fraction of the maximum shared score} achieved, and the \emph{deviation from
optimal} (all-green assignment).  A secondary observable is the \emph{individual-vs-team
trade-off}: participants who choose red more often earn more personally but contribute
less collectively.

\subsection{Repeated Attempts}

Each experimental session consists of up to $A$ attempts (default $A = 3$) on the same
graph instance.  After completing an attempt, the participant sees a results screen
showing the final colouring, a score breakdown, and (from the second attempt onwards)
a comparison with previous attempts.  They may then choose to start another attempt or
end the session.

Each new attempt begins with the graph in the same initial (uncoloured) state; the
previous attempt's colouring is displayed as a visual reference but does not constrain
the new attempt.  In Mode~2, a new plan is generated for each attempt using the
same random seed incremented by the attempt number, so the plan for attempt~$k$ is a
deterministic function of the graph, the seed, and $k$.  Across attempts, planning
quality is held constant (same mode, same $k$-dependent seed), but the sequential
nature of the task means participants who learned from the previous attempt can apply
that knowledge.

%% ============================================================
\section{Graph Topologies}
\label{sec:graphs}
%% ============================================================

The study includes a family of systematically designed graph instances spanning three
cluster sizes and two topological families.  All graphs are defined with nodes partitioned
into two or three agent clusters plus one human cluster, and all use the same colour domain
$\mathcal{C} = \{\text{red}, \text{blue}, \text{green}\}$.

\subsection{Standard Cluster Graphs}

Standard graphs are characterised by dense intra-cluster connectivity and sparse
inter-cluster edges.  Table~\ref{tab:standard_graphs} summarises the variants.

\begin{table}[h]
\centering
\caption{Standard cluster graph variants.  $n$ = total nodes;
         $n_c$ = nodes per cluster; $|E|$ = total edges.}
\label{tab:standard_graphs}
\begin{tabular}{llccc}
\toprule
\textbf{Name} & \textbf{Description} & $n$ & $n_c$ & $|E|$ \\
\midrule
\texttt{s9\_3c}      & 3 clusters $\times$ 3 nodes; internal triangles; H--A and H--B cross-edges  & 9  & 3 & 12 \\
\texttt{s9\_3c\_xab} & \texttt{s9\_3c} + explicit A--B cross-links (2 edges)                      & 9  & 3 & 14 \\
\texttt{s12\_3c}     & 3 clusters $\times$ 4 nodes; quad + diagonal per cluster; 6 cross-edges   & 12 & 4 & 17 \\
\texttt{s12\_3c\_xab}& \texttt{s12\_3c} + A--B cross-links (2 edges)                             & 12 & 4 & 19 \\
\texttt{s18\_3c}     & 3 clusters $\times$ 6 nodes; 6-cycle + 2 chords; 6 cross-edges            & 18 & 6 & 26 \\
\texttt{s18\_3c\_xab}& \texttt{s18\_3c} + A--B cross-links (3 edges)                             & 18 & 6 & 29 \\
\midrule
\texttt{s9\_3c\_c}   & 4-player variant of \texttt{s9\_3c} (adds \AgentC{} cluster)              & 12 & 3 & -- \\
\texttt{s12\_3c\_c}  & 4-player variant of \texttt{s12\_3c}                                       & 16 & 4 & -- \\
\texttt{s18\_3c\_c}  & 4-player variant of \texttt{s18\_3c}                                       & 24 & 6 & -- \\
\bottomrule
\end{tabular}
\end{table}

The \texttt{xab} variants add direct cross-edges between the two agent clusters (not
mediated by the human cluster), increasing inter-agent coupling and reducing the human's
structural centrality in the coordination problem.

\subsection{Hub Topology Graphs}

Hub graphs use a \emph{star topology} within each cluster: one hub node is connected to
all satellite nodes in the cluster.  This topology concentrates constraint propagation
through the hub node, creating a cleaner path-dependency structure.  Table~\ref{tab:hub_graphs}
summarises the hub variants.

\begin{table}[h]
\centering
\caption{Hub topology graph variants.  $n$ = total nodes; architecture is
         the intra-cluster structure; cross-edges is the inter-cluster connection pattern.}
\label{tab:hub_graphs}
\begin{tabular}{llcc}
\toprule
\textbf{Name} & \textbf{Architecture} & $n$ & \textbf{Key cross-edges} \\
\midrule
\texttt{hub12}      & Star (hub + 3 satellites) per cluster & 12 & $h_1$--$a_1$, $h_1$--$b_1$ \\
\texttt{hub12\_tri} & \texttt{hub12} + $a_1$--$b_1$                                  & 12 & Hub triangle \\
\texttt{hub15\_deep}& Hub triangle + $h_2$--$a_1$, $h_3$--$b_1$                    & 15 & Deep cross-constraints \\
\texttt{hub18}      & Hub + satellite pairs per cluster; full hub triangle           & 18 & Hub triangle + pair threading \\
\texttt{hub18b}     & Asymmetric: H chain, A disjoint pairs, B chain                & 18 & Hub triangle + asymmetric links \\
\bottomrule
\end{tabular}
\end{table}

Hub topologies have the property that the hub node's colour forces or strongly
constrains the colours of adjacent cluster hubs (via inter-cluster edges).  In
\texttt{hub12}, for example, colouring $h_1$ (the human hub) to red immediately
constrains $a_1$ (the \AgentA{} hub) to be non-red, and $b_1$ (the \AgentB{} hub) to
also be non-red.  This cascades through both agent clusters.

The asymmetric \texttt{hub18b} variant models realistic organisational heterogeneity:
the human cluster and \AgentB{} cluster have chain internal structure (which propagates
colour constraints along the chain), while \AgentA{}'s cluster has disjoint pairs
(which do not propagate constraints as strongly).  This creates an interesting
scenario where the human's constraints ripple through structurally different agent
clusters at different rates.

\subsection{Colouring Order}

In the default (non-random) turn order, nodes are coloured in a fixed sequence:
all human-cluster nodes, then all \AgentA-cluster nodes, then all \AgentB-cluster nodes.
Within each cluster, nodes follow the order defined in the graph specification (typically
hub first, then satellites in index order).

This ordering ensures that the human acts first and sets the context for both agents.
The agents' automatic colouring decisions (in Mode~1) or plan steps (in Mode~2) are
then presented after the human has committed their choices, reflecting the causal
structure of the task.

%% ============================================================
\section{Experimental Conditions}
\label{sec:conditions}
%% ============================================================

The study uses a within-subjects design with three conditions applied to each participant.
The conditions are:

\begin{description}[leftmargin=2em, itemsep=4pt]
  \item[Mode~1 — Manual Collaborative]
    Per-node proposals.  For each human-owned node, both agents
    simultaneously issue a proposed colour accompanied by a written rationale.
    A score-prediction table shows each agent's expected outcome under each
    possible colour choice.  The human reads these proposals and then freely
    chooses any colour.  After all human nodes are coloured, agent nodes are
    coloured automatically by the agents' own greedy logic.

  \item[Mode~2A — Autonomous, High Quality]
    The planning agent generates a complete proposed colouring for all nodes
    before the session starts, using a quality parameter of $q = 0.8$: each
    plan step independently picks the shared-optimal colour with probability
    0.8, and picks a random safe colour with probability 0.2.  This plan is
    presented to the human one step at a time.  For human-owned nodes, the
    participant can accept the plan's recommendation or override it.  For
    agent-owned nodes, the plan step is applied automatically with a brief
    visual pause.

  \item[Mode~2B — Autonomous, Low Quality]
    Identical in interface and interaction structure to Mode~2A.  The
    planning agent uses quality parameter $q = 0.3$: each plan step picks
    the shared-optimal colour with probability 0.3 only, making the plan
    substantially worse on average.  Participants are not informed of the
    quality level; the interface is intentionally indistinguishable from
    Mode~2A.
\end{description}

\paragraph{Within-subjects counterbalancing.}
All participants complete all three conditions.  Conditions are presented in
counterbalanced order across participants to control for learning and fatigue effects.
A different graph instance (from the families described in Section~\ref{sec:graphs})
is used for each condition.  Graph instances are matched on the number of nodes,
cluster sizes, and approximate scoring difficulty.

\paragraph{Design rationale.}
Mode~1 represents a \emph{transparent advisory} interface: the agent is fully
explainable at the per-node level and the human always retains final decision authority.
Mode~2A represents a \emph{competent autonomous} system: the agent plans well and the
human's role shifts toward review and approval.  Mode~2B represents an
\emph{incompetent autonomous} system that \emph{looks} identical to Mode~2A.  The
comparison between Mode~2A and Mode~2B is the primary test of whether participants
appropriately calibrate trust based on the quality of outcomes rather than the form of
the interface.  The comparison between Mode~1 and Mode~2 (pooling 2A and 2B) tests
the effect of interaction style: per-node deliberation versus pre-planned review.

%% ============================================================
\section{Agent Models}
\label{sec:agents}
%% ============================================================

All agent logic is deterministic and rule-based.  No large language model or learned
model is used.  This design choice ensures full reproducibility, eliminates variance due
to stochastic model outputs, and allows precise control over agent quality.  The term
``agent'' is used throughout the interface; participants are not told the mechanism.

\subsection{Proposal Agent (Mode~1)}
\label{sec:proposal_agent}

The proposal agent (\texttt{ProposalAgent}) generates a per-node colour recommendation
each time a new human-owned node becomes active.  Each AI agent runs this logic
independently; both proposals are shown simultaneously.

\paragraph{Adjacency analysis.}
The agent first identifies which of its own nodes are adjacent (share an edge) to the
current human node.  If the agent has no adjacent node, it does not generate a proposal
for this step.  If multiple adjacent nodes exist, the agent selects one reference node
using a priority rule: already-coloured adjacent nodes take precedence (actionable
advice); among uncoloured nodes, the first in alphabetical order is chosen.

\paragraph{Colour recommendation.}
The recommendation logic depends on the reference node's status:
\begin{itemize}[nosep]
  \item \textbf{Reference node already coloured:} Recommend the first colour (in
    $\{\text{red}, \text{blue}, \text{green}\}$ order) that does not clash with the
    reference node's colour.  This avoids a cross-cluster penalty.
  \item \textbf{Reference node not yet coloured:} The agent anticipates placing its own
    preferred colour on the reference node later, so it recommends a colour for the
    human node that avoids clashing with that anticipated choice.  Specifically, it
    recommends the first colour that is \emph{not} the agent's preferred colour.
\end{itemize}
If no non-clashing colour exists (e.g.\ all colours are ruled out simultaneously), the
agent falls back to the first available colour in domain order.

\paragraph{Rationale generation.}
Each proposal includes a short rationale (one or two sentences displayed inline) and a
full rationale (expandable detail, shown on user request).  Short rationales follow one
of two templates:

\begin{itemize}[nosep]
  \item \textit{Clash avoidance}: ``My node $X$ is already [colour] --- I recommend
    [colour] here so our adjacent nodes don't clash.''
  \item \textit{Future planning}: ``My node $X$ is next to this one and I plan to colour
    it [preferred] --- so I recommend [colour] here to keep our nodes well-coordinated.''
\end{itemize}

When the attempt number is greater than 1, a note is added comparing the current
recommendation to the previous attempt's colour at this node: ``Last attempt this was
[colour]; my recommendation is [colour] this time.''

Full rationales include: the adjacent node identification, the rationale for colour
selection, a per-colour point table for all participants, shared-score contributions per
colour, and a count of remaining uncoloured nodes.

\paragraph{Expected score display.}
In addition to the written rationale, Mode~1 shows a score-prediction table (described
in Section~\ref{sec:interface}) showing each agent's expected score contribution under
each possible colour choice for the current node.  This is computed using a forward
simulation: for each possible colour $c$ that the human might choose, the agent
greedily colours its remaining nodes (preferring clash avoidance, then score
maximisation), and reports the agent's resulting total score.  This gives participants
explicit quantitative information about the consequences of their choices for each
agent.

\paragraph{Self-colouring (agent-owned nodes).}
When an agent-owned node becomes active (after all human nodes are decided), the agent
colours it automatically using a greedy policy: avoid clashing with already-coloured
adjacent nodes; among clash-free candidates, pick the colour that maximises the agent's
own score for that node.

\subsection{Planning Agent (Mode~2)}
\label{sec:planning_agent}

The planning agent (\texttt{PlanningAgent}) produces a complete proposed colouring
for all nodes \emph{before} the session begins.  This plan is then delivered to the
participant step by step.

\paragraph{Deliberation process.}
The planning agent is initialised with a quality parameter $q \in [0, 1]$.  In the
study, $q = 0.8$ for Mode~2A and $q = 0.3$ for Mode~2B.  A seeded random number
generator is initialised with seed $s + k$, where $s$ is the session seed and $k$ is
the attempt number (1-based), ensuring different but reproducible plans across attempts.

For each node in the colouring sequence, the deliberation process runs as follows:

\begin{enumerate}[nosep]
  \item \textbf{Clash detection:} Identify all colours already assigned to adjacent
    nodes (nodes already coloured in the current plan-in-progress).  Colours that
    would clash with an adjacent already-planned colour are excluded; the remaining
    colours are the \emph{safe candidates}.  If all colours clash (which can occur in
    dense graphs), all colours are re-admitted as candidates.

  \item \textbf{Agent votes:} Each sub-agent independently ``votes'' for its preferred
    colour (the colour that maximises its own per-node points, regardless of the graph
    context).  Votes are logged in the deliberation record.

  \item \textbf{Quality-biased selection:} A random draw determines the selection
    mechanism for this step:
    \begin{itemize}[nosep]
      \item With probability $q$: select the safe candidate with the highest
        shared-score total (the shared-optimal safe colour).  This is the
        ``high-quality'' choice.
      \item With probability $1 - q$: select uniformly at random from the safe
        candidates.  This may coincidentally be the shared-optimal colour, but is not
        guided by shared score.
    \end{itemize}

  \item \textbf{Rationale generation:} A short rationale and full rationale are
    generated, templated as:
    \begin{itemize}[nosep]
      \item Short (shared-optimal): ``[Colour] gives the team $N$ pts on this node,
        which maximises our shared score.''
      \item Short (random): ``[Colour] was selected for [node] after agent deliberation
        (team score: $N$ pts).''
    \end{itemize}
    The full rationale includes all agent votes, the vote tally per colour, the
    shared-score contribution of each colour, the selection method, and the final
    choice.
\end{enumerate}

\paragraph{Deliberation log.}
The complete deliberation log is accessible to participants via a collapsible panel in
the interface (labelled ``How was this plan made?'').  This log shows the step-by-step
deliberation for every node: which colours were safe, how each agent voted, which
selection method was used, and what the expected score contribution is.  An interaction
event is logged whenever a participant opens this panel.

\paragraph{Expected score.}
The expected shared score for the full plan is computed as the sum of the shared-score
contributions of the chosen colour at each node, before considering any clashes.  This
figure is displayed in the plan summary panel.

\paragraph{Plan invariance.}
Because the planning process uses a seeded RNG, the same graph instance, seed, and
attempt number always produce the same plan.  This means Mode~2A and Mode~2B plans are
reproducible and their quality is deterministic.  Within a condition, different
participants see different plans only if they use different seeds (configurable at
session launch).

%% ============================================================
\section{Interface}
\label{sec:interface}
%% ============================================================

The experiment interface is implemented in Python using the Tkinter GUI toolkit.  All
conditions use the same application window and layout; only the right-hand panel changes
between Mode~1 and Mode~2.

\subsection{Window Layout}

The main window ($1020 \times 660$ minimum) is divided into three horizontal zones:

\begin{description}[leftmargin=2em, itemsep=2pt]
  \item[Top — Scoring HUD]  A horizontal bar showing the current shared score (large,
    green) and individual score cards for each participant (colour chip + name + points).
    Updates after every colour decision.
  \item[Middle — Main Panel]  Divided into a left pane (graph display) and a right
    pane (mode-specific interaction panel).
  \item[Bottom — Sequence Bar]  A breadcrumb row of small cells, one per node,
    showing colouring progress.  Grey cells are uncoloured; filled cells show the
    assigned colour.
\end{description}

\subsection{Graph Display (Left Pane)}

The left pane contains two graph canvases stacked vertically:

\begin{itemize}[nosep]
  \item \textbf{Overview canvas} ($440 \times 280$~px): Shows the complete graph with
    all nodes and edges.  Nodes are positioned according to precomputed fractional
    coordinates stored per graph preset.  Coloured nodes show their assigned colour;
    uncoloured nodes are grey.
  \item \textbf{Neighbourhood canvas} ($440 \times 210$~px): Shows the current active
    node and its direct neighbours in isolation.  Provides a local view that highlights
    the immediate constraint context.
\end{itemize}

\subsection{Mode~1 Right Panel}

The Mode~1 panel contains the following elements from top to bottom:

\begin{enumerate}[nosep]
  \item \textbf{Reward key}: A small table showing the point values of each colour for
    the human participant, with the highest-value colour highlighted.
  \item \textbf{Current node label}: Displays the name of the active node.
  \item \textbf{Score prediction table}: A grid with rows for each agent that has an
    adjacent node, and columns for each colour.  Each cell shows the agent's predicted
    score contribution if the human chooses that colour, with a star ($\star$) marking
    the agent's own preferred colour.
  \item \textbf{Colour choice buttons}: One button per colour, enabled only after
    proposals have been shown.  Clicking a button records the decision and advances to
    the next node.
\end{enumerate}

Agent proposals are shown automatically when a human node becomes active.  Both agents'
proposals and the score table appear simultaneously.  When the human clicks a colour
button, the decision is logged (including source = ``human'', decision time, and whether
a majority of agents agreed with the choice).

\subsection{Mode~2 Right Panel}

The Mode~2 panel presents the current plan step:

\begin{enumerate}[nosep]
  \item \textbf{Proposed colour}: The plan's recommended colour for the current node,
    displayed prominently.
  \item \textbf{Short rationale}: One or two sentences from the planning agent's
    rationale for this recommendation, displayed below the proposed colour.
  \item \textbf{Action buttons}: Three buttons --- \textit{Accept} (records the plan's
    colour, source = ``accepted\_plan''), \textit{Explain} (toggles the full rationale
    text), and \textit{Modify} (opens a sub-panel with colour buttons for overriding
    the plan, source = ``modified\_plan'').
  \item \textbf{Full rationale panel}: A scrollable text field (hidden by default,
    toggled by Explain) showing the complete rationale for this step including all agent
    votes, vote tallies, and shared-score contributions per colour.
  \item \textbf{Deliberation log panel}: A collapsible text field accessible via a
    labelled checkbox, showing the complete planning deliberation log for the whole
    session.
\end{enumerate}

For agent-owned nodes in Mode~2, the plan step is applied automatically with a 500~ms
visual pause, and no buttons are shown.  The participant observes the plan being executed
but cannot intervene.

\subsection{End-of-Attempt Results Panel}

When all nodes have been coloured, the colouring panel is replaced in-place by a results
view showing:

\begin{itemize}[nosep]
  \item A final graph rendering with the completed colouring.
  \item A score table: each participant's name, preferred colour, individual score, and
    clash penalty (if any).
  \item The \textbf{shared total} score, displayed prominently.
  \item A score history table (from the second attempt onwards) listing the shared
    total for each attempt, with markers indicating the current attempt and the best
    attempt.
  \item \textbf{Action buttons}: ``Start Next Attempt'' (if attempts remain) and
    ``Finish Session''.
\end{itemize}

The previous attempt's colouring is preloaded as a visual overlay on the graph at the
start of each new attempt (with a 200~ms delay for canvas layout), giving participants a
reference without constraining their choices.

\subsection{Event Logging}

All interactions are written to a JSONL event log at\\
\texttt{results/participants/<pid>\_<timestamp>/events.jsonl}.  Each line is a JSON
object with a Unix timestamp and event type.  Table~\ref{tab:events} lists the key
logged events.

\begin{table}[h]
\centering
\caption{Key JSONL events logged during each session.
         All events include a Unix timestamp \texttt{ts} and event type \texttt{event}.}
\label{tab:events}
\begin{tabularx}{\linewidth}{lX}
\toprule
\textbf{Event type}          & \textbf{Fields logged} \\
\midrule
\texttt{session\_start}      & participant\_id, mode, graph, colours, seed, max\_attempts, points\_config \\
\texttt{attempt\_start}      & attempt number (1-based) \\
\texttt{node\_ready}         & attempt, node, node\_index (starts per-node timer) \\
\texttt{proposals\_shown}    & (Mode~1) attempt, node, proposals [\{agent, colour\}] \\
\texttt{explanation\_click}  & (Mode~1) attempt, node, agent name \\
\texttt{colour\_chosen}      & attempt, node, colour, source, decision\_time\_s, agents\_agreed \\
\texttt{plan\_step\_shown}   & (Mode~2) attempt, node, plan\_colour \\
\texttt{plan\_accepted}      & (Mode~2) attempt, node, plan\_colour \\
\texttt{plan\_modified}      & (Mode~2) attempt, node, plan\_colour, chosen\_colour \\
\texttt{plan\_explanation\_clicked} & (Mode~2) attempt, node \\
\texttt{deliberation\_log\_opened}  & (Mode~2) attempt \\
\texttt{attempt\_complete}   & attempt, assignment \{node: colour\}, score \{total, by\_owner\}, elapsed\_s \\
\texttt{replay\_requested}   & attempt (user chose to start another attempt) \\
\texttt{session\_end}        & total\_attempts, score\_trajectory [int], session\_elapsed\_s \\
\bottomrule
\end{tabularx}
\end{table}

The \texttt{colour\_chosen} event includes \texttt{source} $\in$ \{``human'',
``accepted\_plan'', ``modified\_plan'', ``agent\_auto''\}, distinguishing human
initiative, plan acceptance, plan override, and automatic agent colouring.  The
\texttt{agents\_agreed} field (boolean) records whether a majority of agents whose
proposals or plan recommended this node had recommended the chosen colour.
Decision time is measured from the moment the node was made active
(\texttt{node\_ready}) to the moment the colour button was pressed.

%% ============================================================
\section{Procedure}
\label{sec:procedure}
%% ============================================================

\subsection{Session Structure}

Each session is launched via a graphical configuration screen where the experimenter (or
participant, in unattended conditions) enters the participant ID, selects the mode and
graph preset, sets the random seed, and specifies the maximum number of attempts.  A
subprocess then opens the experiment window.

A session consists of the following stages:

\begin{enumerate}[nosep]
  \item \textbf{Introduction and practice}: Participants are introduced to the graph
    colouring task, the scoring system, and the interface.  (Specific instructions TBD
    per study protocol.)
  \item \textbf{Experimental task}: The participant completes the colouring task in the
    assigned condition, with up to $A = 3$ attempts.
  \item \textbf{Post-condition questionnaire}: After each condition, participants
    complete a questionnaire on workload, trust, and perceived agent quality.  (Scale
    details TBD.)
\end{enumerate}

\subsection{Within-Attempt Flow}

Within each attempt, the interaction proceeds as follows:

\begin{enumerate}[nosep]
  \item For Mode~2 only: the planning agent generates the complete plan (logged
    internally; plan generation is not visible to the participant).
  \item The first node in the sequence is activated.
  \item For each human-owned node:
    \begin{itemize}[nosep]
      \item Mode~1: Agent proposals appear; the participant reads and chooses a colour.
      \item Mode~2: The plan's recommendation appears; the participant accepts or overrides.
    \end{itemize}
  \item For each agent-owned node: the agent's colour is applied automatically (Mode~1:
    greedy rule; Mode~2: plan step), with a 500~ms visual pause.
  \item When all nodes are coloured, the results panel appears.  The score is displayed
    and compared with previous attempts.
  \item The participant chooses ``Start Next Attempt'' or ``Finish Session''.
\end{enumerate}

\subsection{Randomisation and Reproducibility}

All stochastic elements (plan generation, optional node order shuffle) are controlled by
a session seed $s$ (default: 42).  The plan for attempt $k$ uses RNG seed $s + k$.
If random turn order is enabled, the node sequence for attempt $k$ is derived from
$s + k$ before plan generation, so plan and presentation order are always consistent.
This ensures full reproducibility: given the same participant ID, mode, graph, and seed,
the entire session is deterministic.

%% ============================================================
\section{Dependent Variables and Hypotheses}
\label{sec:measures}
%% ============================================================

\subsection{Primary Outcome: Shared Score}

The primary dependent variable is the \emph{shared score} achieved at the end of each
attempt.  This is the sum of all participants' individual scores after applying clash
penalties (Equation~\ref{eq:score}).  The maximum achievable shared score is $8|V|$
(all nodes green, no clashes).  We report both raw scores and \emph{efficiency}:
$\eta = \text{Score}_\Shared / (8|V|)$.

\subsection{Secondary Outcomes}

\paragraph{Individual vs.\ shared trade-off.}
The human participant's individual score relative to the shared score.  A participant who
consistently chooses red (their preferred colour) earns higher individual scores but lower
shared scores.

\paragraph{Plan modification rate (Mode~2 only).}
The proportion of human-node plan steps that the participant overrides
(\texttt{source} = ``modified\_plan'').  High modification rate in Mode~2B (low quality)
relative to Mode~2A (high quality) is evidence of appropriate trust calibration; similar
modification rates across 2A and 2B would suggest participants cannot detect quality
differences.

\paragraph{Explanation engagement.}
The proportion of nodes for which the participant expanded the full rationale (Mode~1:
\texttt{explanation\_click}; Mode~2: \texttt{plan\_explanation\_clicked}).  Also, the
proportion of attempts in which the deliberation log was opened
(\texttt{deliberation\_log\_opened}).

\paragraph{Decision latency.}
The mean \texttt{decision\_time\_s} per node, computed from event timestamps.  Longer
decision times may indicate deeper deliberation; shorter times may indicate automatic
acceptance.

\paragraph{Score improvement across attempts.}
The difference in shared score between attempt $k$ and attempt 1, $\Delta_k =
\text{Score}^{(k)} - \text{Score}^{(1)}$.  Positive $\Delta$ indicates learning;
negative $\Delta$ would be unexpected and may indicate fatigue or strategy abandonment.

\paragraph{Agents-agreed rate.}
The proportion of decisions for which the majority of agents recommended the colour the
participant ultimately chose (\texttt{agents\_agreed} = True).  In Mode~1 this reflects
whether the participant followed agent advice; in Mode~2 it is equivalent to the
plan-acceptance rate.

\subsection{Hypotheses}

\begin{description}[leftmargin=2em, itemsep=3pt]
  \item[H1 \textit{(Quality detection)}]
    Participants will modify a higher proportion of plan steps in Mode~2B than in
    Mode~2A, showing appropriate behavioural trust calibration based on perceived plan
    quality.

  \item[H2 \textit{(Score gap)}]
    Mode~2A will produce higher shared scores than Mode~2B, and the gap will be larger
    when modification rates are well-calibrated (i.e., participants override low-quality
    steps without over-riding high-quality ones).

  \item[H3 \textit{(Transparency benefit)}]
    Mode~1 will produce better-calibrated decisions than Mode~2, because per-node
    proposals with explicit rationales give participants more material to evaluate each
    recommendation individually, compared to a pre-committed plan.

  \item[H4 \textit{(Explanation engagement and calibration)}]
    Participants who engage with full rationales (explanation click rate above the
    median) will show better plan modification calibration in Mode~2: they will override
    low-quality plan steps more often while accepting high-quality steps at similar rates.

  \item[H5 \textit{(Learning across attempts)}]
    Shared scores will improve significantly from attempt~1 to attempt~3, with the
    improvement larger in Mode~2B than Mode~2A (because there is more room to improve by
    overriding bad recommendations) and larger in Mode~2 than Mode~1 (because learning
    the agent's plan tendencies is more useful when the same plan structure recurs).

  \item[H6 \textit{(Individual vs.\ shared sacrifice)}]
    Participants who accept more agent recommendations (higher agents-agreed rate) will
    show lower individual scores but higher shared scores, consistent with trading
    personal gain for collective benefit as the agents advise.
\end{description}

%% ============================================================
\section{Results}
\label{sec:results}
%% ============================================================

\textit{This section will be completed once data collection is finished.  The subsections
below provide the planned analysis structure.}

\subsection{Overview of Data Collected}

[Table: participant count, conditions completed, attempt counts, any exclusions.]

\subsection{Primary Outcome: Shared Score by Condition}

[ANOVA or mixed-effects model on shared score with condition (Mode~1, 2A, 2B) as the
within-subjects factor and attempt number as a covariate.  Report means, standard
deviations, and effect sizes.  Include a figure showing mean shared score per condition
across attempts.]

\subsection{Plan Modification Rate (Modes~2A vs 2B)}

[Paired comparison of modification rates between Mode~2A and Mode~2B.  Report whether
participants modify more steps in Mode~2B (H1).  Break down by node position in sequence
to test whether modification rates change as participants progress through an attempt.]

\subsection{Explanation Engagement}

[Descriptive statistics of explanation click rates per condition.  Test H4: do
high-engagers show better calibration?  Report median split analysis and correlation
between engagement and plan modification calibration.]

\subsection{Decision Latency}

[Mean decision time per node by condition, with standard error.  Compare latency for
accepted-plan vs.\ modified-plan steps (Mode~2) and for conditions with vs.\ without
explanation clicks (Mode~1).]

\subsection{Score Improvement Across Attempts}

[Plot score trajectory (score vs.\ attempt number) per condition.  Test H5 with a
linear trend test on score across attempts.  Report improvement magnitude $\Delta_3
= \text{Score}^{(3)} - \text{Score}^{(1)}$ by condition.]

\subsection{Individual vs.\ Shared Trade-off}

[Scatter plot: individual score vs.\ shared score per participant per condition.  Report
correlation between agents-agreed rate and (a) individual score and (b) shared score,
testing H6.]

%% ============================================================
\section{Discussion}
\label{sec:discussion}
%% ============================================================

\textit{To be completed.  Planned discussion points:}

\begin{itemize}[nosep]
  \item Whether the interface indistinguishability of Mode~2A and Mode~2B produces
    overtrust (low modification in 2B) or whether participants can detect quality
    through outcome patterns across nodes.
  \item The role of rationale transparency: does Mode~1's per-node explainability
    produce meaningfully different behaviour from Mode~2's global plan transparency?
  \item Path dependency effects: are early-sequence decisions predictive of final
    shared score, and do participants who make globally better early choices (e.g.,
    choosing green for high-adjacency nodes) achieve higher totals?
  \item Comparison with a na\"{i}ve baseline (all human nodes red, all agents optimal)
    to contextualise score achievements.
  \item Implications for the design of autonomous AI systems: what transparency
    mechanisms are necessary for users to detect and compensate for poor-quality
    autonomous plans?
\end{itemize}

%% ============================================================
\section{Conclusion}
\label{sec:conclusion}
%% ============================================================

\textit{To be completed.}

We have described a controlled experimental platform for studying trust calibration and
collaborative decision-making in sequential human--agent graph colouring.  The platform's
key features are: (1) a payoff structure that creates an intrinsic tension between
individual and collective optimisation; (2) a sequential colouring order that introduces
path dependency; (3) three experimental conditions that vary agent autonomy and plan
quality while holding the interface constant; and (4) fine-grained event logging that
supports analysis of decision timing, explanation engagement, and score trajectories
across repeated attempts.

The distinguishing feature of the design is that Mode~2A and Mode~2B are
\emph{superficially identical}: same visual layout, same language, same plan-presentation
mechanism.  Whether participants detect the quality difference and calibrate their
override behaviour accordingly is the central empirical question.  A finding of overtrust
(no calibration) would suggest that interface presentation alone determines reliance
patterns, regardless of actual quality.  A finding of appropriate calibration would
suggest that participants can perform implicit quality inference from outcome patterns
within a single session.

\bibliographystyle{plainnat}
\bibliography{refs}

\clearpage
\appendix
\renewcommand{\thesection}{\Alph{section}}

%% ============================================================
\section{Scoring Model: Formal Specification}
\label{app:scoring}
%% ============================================================

\subsection{Point Configuration}

Let $\mathcal{P} = \{h, a_1, a_2\}$ be the set of participants.  The point function
$\text{pts} : \mathcal{P} \times \mathcal{C} \to \mathbb{Z}_{\geq 0}$ is given by
Table~\ref{tab:default_points} in the main text.  The shared-score contribution of
colour $c$ (regardless of node or participant) is:

\[
  \text{shared}(c) = \sum_{p \in \mathcal{P}} \text{pts}(p, c)
\]

In the default configuration: $\text{shared}(\text{red}) = 7$,
$\text{shared}(\text{blue}) = 7$, $\text{shared}(\text{green}) = 8$.

\subsection{Clash Penalty Allocation}

For edge $(u, v) \in E$ with $x_u = x_v$ (same colour), the clash penalty of
$\delta = 10$ is allocated as follows.  Let $\mathrm{owner}(u)$ denote the participant
who owns node $u$.

\[
  \text{penalty}(u, v, p) = \begin{cases}
    \delta & \text{if } \mathrm{owner}(u) = \mathrm{owner}(v) = p \\
    \lfloor \delta / 2 \rfloor + (\delta \bmod 2) & \text{if } \mathrm{owner}(u) = p \neq \mathrm{owner}(v) \\
    \lfloor \delta / 2 \rfloor & \text{if } \mathrm{owner}(v) = p \neq \mathrm{owner}(u) \\
    0 & \text{otherwise}
  \end{cases}
\]

For $\delta = 10$: intra-participant clashes cost 10 to the owning participant;
inter-participant clashes cost 5 to each owner (since $10 \bmod 2 = 0$).

\subsection{Optimal and Baseline Scores}

For a graph with $|V| = n$ nodes and $|E_{\times}|$ inter-cluster edges:

\begin{itemize}[nosep]
  \item \textbf{Upper bound (no clashes, all green)}: $\text{shared}(\text{green}) \times n = 8n$.
  \item \textbf{Penalty for all-red}:  $\text{shared}(\text{red}) \times n = 7n$.
  \item \textbf{Penalty if clashes occur}: each clash deducts 10 from the shared total.
\end{itemize}

The difference between all-green and all-red is $n$ shared points.  For the default
12-node graph, this is 12 shared points, a 12.5\% difference on the theoretical maximum.

%% ============================================================
\section{Planning Agent: Algorithm}
\label{app:planning}
%% ============================================================

Algorithm~\ref{alg:plan} gives the pseudocode for the planning agent's deliberation
process.

\begin{algorithm}
\caption{PlanningAgent.generate\_plan($k$, prior\_attempts, node\_order)}
\label{alg:plan}
\begin{algorithmic}[1]
\REQUIRE Attempt number $k$, list of prior assignments, node sequence $\sigma$,
         quality $q \in [0,1]$, seed $s$
\STATE Initialise $\mathrm{rng} \gets \text{Random}(s + k)$
\STATE $\mathrm{assignment} \gets \{\}$; $\mathrm{plan} \gets []$
\FORALL{$v$ in $\sigma$}
  \STATE $\mathrm{adj\_colours} \gets \{x_u \mid (u,v) \in E, u \in \mathrm{assignment}\}$
  \STATE $\mathrm{candidates} \gets \mathcal{C} \setminus \mathrm{adj\_colours}$
  \IF{$\mathrm{candidates} = \emptyset$}
    \STATE $\mathrm{candidates} \gets \mathcal{C}$ \COMMENT{admit all if full clash}
  \ENDIF
  \STATE \texttt{// Sub-agent votes (self-interested, ignored for selection)}
  \FORALL{agent $a$}
    \STATE vote $\gets \argmax_{c} \text{pts}(a, c)$
  \ENDFOR
  \STATE $c^* \gets \argmax_{c \in \mathrm{candidates}} \text{shared}(c)$ \COMMENT{best safe colour}
  \IF{$\mathrm{rng.random()} < q$}
    \STATE $\mathrm{chosen} \gets c^*$ \COMMENT{quality: pick shared-optimal}
  \ELSE
    \STATE $\mathrm{chosen} \gets \mathrm{rng.choice}(\mathrm{candidates})$ \COMMENT{explore randomly}
  \ENDIF
  \STATE $\mathrm{assignment}[v] \gets \mathrm{chosen}$
  \STATE Append $\mathrm{PlannedStep}(v, \mathrm{chosen}, \ldots)$ to $\mathrm{plan}$
\ENDFOR
\RETURN $\mathrm{Plan}(\mathrm{plan}, \sum_v \text{shared}(\mathrm{assignment}[v]), k)$
\end{algorithmic}
\end{algorithm}

Note that sub-agent votes are recorded in the deliberation log (and shown to participants
in the full rationale) but do not affect the colour selection.  The selection mechanism
depends solely on the quality draw and the shared-score ordering of safe candidates.

%% ============================================================
\section{Proposal Agent: Algorithm}
\label{app:proposal}
%% ============================================================

Algorithm~\ref{alg:proposal} gives the pseudocode for the proposal agent's per-node
recommendation.

\begin{algorithm}
\caption{ProposalAgent.propose\_for\_node($v$, assignment\_so\_far, attempt)}
\label{alg:proposal}
\begin{algorithmic}[1]
\REQUIRE Current node $v$, partial assignment so far, attempt number
\STATE $\mathrm{adjacent} \gets \{u \in V_a \mid (u,v) \in E\}$
\IF{$\mathrm{adjacent} = \emptyset$}
  \RETURN None \COMMENT{no proposal if no adjacent agent node}
\ENDIF
\STATE $\mathrm{coloured\_adj} \gets \{u \in \mathrm{adjacent} \mid u \in \mathrm{assignment\_so\_far}\}$
\IF{$\mathrm{coloured\_adj} \neq \emptyset$}
  \STATE $\mathrm{ref} \gets$ any $u \in \mathrm{coloured\_adj}$
  \STATE $\mathrm{avoid} \gets \{\mathrm{assignment\_so\_far}[\mathrm{ref}]\}$
  \STATE $\mathrm{rationale\_type} \gets$ ``clash\_avoidance''
\ELSE
  \STATE $\mathrm{ref} \gets$ first of $\mathrm{adjacent}$ alphabetically
  \STATE $\mathrm{avoid} \gets \{\mathrm{preferred\_colour}(a)\}$
  \STATE $\mathrm{rationale\_type} \gets$ ``future\_planning''
\ENDIF
\STATE $\mathrm{candidates} \gets \mathcal{C} \setminus \mathrm{avoid}$
\IF{$\mathrm{candidates} = \emptyset$}
  \STATE $\mathrm{candidates} \gets \mathcal{C}$
\ENDIF
\STATE $\mathrm{recommended} \gets \mathrm{candidates}[0]$ \COMMENT{first in domain order}
\RETURN $\mathrm{NodeProposal}(v, \mathrm{recommended}, \text{rationale}(\mathrm{rationale\_type}, \ldots))$
\end{algorithmic}
\end{algorithm}

%% ============================================================
\section{Event Log Schema}
\label{app:schema}
%% ============================================================

Each line of the JSONL log is a JSON object.  Below is the complete field specification
for the most important events.

\paragraph{\texttt{colour\_chosen}}
\begin{verbatim}
{
  "ts": <float>,            // Unix timestamp
  "event": "colour_chosen",
  "attempt": <int>,         // 1-based
  "node": <str>,            // node identifier (e.g. "H1")
  "colour": <str>,          // "red" | "blue" | "green"
  "source": <str>,          // "human" | "accepted_plan" |
                            //  "modified_plan" | "agent_auto"
  "decision_time_s": <float>, // seconds from node_ready to click
  "agents_agreed": <bool>   // majority of agents recommended this colour
}
\end{verbatim}

\paragraph{\texttt{attempt\_complete}}
\begin{verbatim}
{
  "ts": <float>,
  "event": "attempt_complete",
  "attempt": <int>,
  "assignment": { <node>: <colour>, ... },
  "score": {
    "total": <int>,
    "by_owner": { <participant>: <int>, ... }
  },
  "elapsed_s": <float>
}
\end{verbatim}

\paragraph{\texttt{session\_end}}
\begin{verbatim}
{
  "ts": <float>,
  "event": "session_end",
  "total_attempts": <int>,
  "score_trajectory": [<int>, ...],  // shared score per attempt
  "session_elapsed_s": <float>
}
\end{verbatim}

\end{document}
